\documentclass[10pt,a4paper]{amsart}
\usepackage[margin=19mm]{geometry}
\usepackage{amsmath,amssymb,mathtools}
\usepackage[scr=rsfs,cal=euler]{mathalfa}
\usepackage{xfrac}
\usepackage[unicode,hidelinks,bookmarks=false]{hyperref}
\newcommand{\ie}{i.e.\ }
\newcommand{\cf}{cf.\ }
\newcommand{\resp}{respectively\ }
\newcommand{\Subsection}{\subsection}
\providecommand{\nobreakdash}{\nobreak\hbox{-}}
\setlength{\emergencystretch}{2em}
\multlinegap0pt
\usepackage[shortlabels]{enumitem}
\usepackage{amssymb}
\usepackage{mathtools}
\usepackage{tikz}
\newtheorem{definition}{Definition}
\newtheorem{Hypotheses}{Hypotheses}
\newtheorem{theorem}{Theorem}
\newtheorem{proposition}{Proposition}
\DeclareMathOperator{\Ext}{Ext}
\newcommand{\R}{\mathbb{R}}
\newcommand{\id}{\mathrm{id}}
\title{Two characters on one punctured Riemann surface}
\author{Pradip Kumar}
\date{Source: arXiv:2603.01905v1 (CC BY 4.0)}
\begin{document}
\maketitle

\noindent\emph{Source and licence.} Two characters on one punctured Riemann surface. Version arXiv:2603.01905v1 (CC BY 4.0), distributed by arXiv under the Creative Commons Attribution 4.0 International licence, http://creativecommons.org/licenses/by/4.0/, as recorded in the arXiv licence metadata for this exact version and shown on the abstract page https://arxiv.org/abs/2603.01905v1. Full licence notice: \texttt{licenses/arxiv-2603.01905-CC-BY-4.0.txt}.\par\smallskip

\noindent\textit{Editorial scope.} Counted unit: the article's theorem thm:good\_example (source0.tex lines 1213--1219). The article's complete original proof of it is delivered with it (source0.tex lines 1221--1288, one \texttt{proof} environment of 68 lines). No mathematics is written, repaired or re-derived: the statement and the proof are the article's own lines, in the article's own notation, and no retained line is substituted or re-flowed.  Retained uncounted as the source-local context the proof needs: lines 301--309; lines 313--340; lines 344--350; lines 560--572; lines 592--604; lines 613--620; lines 622--647; lines 1126--1137; lines 1182--1201.  Nothing was omitted from any retained span, so the package carries no omission record.  No retained line contains a citation, so the wrapper carries no bibliography and no bibliography entry.  No retained line contains a figure environment, an image-inclusion command or drawing code, so no image is delivered and the package declares no asset dependency.  Editorial formatting only, in addition: an AMS \texttt{amsart} preamble that restates the article's own declarations and macros used by the retained lines, namely the environments \texttt{definition}, \texttt{Hypotheses}, \texttt{theorem}, \texttt{proposition} on separate counters, together with the standard packages those lines need.  The retained environments print in this document as Definition 1, Hypotheses 1, Theorem 1, Proposition 1 in order of appearance, whereas the article prints the same items with the numbering of its own full document.  The complete native source of the article as distributed in the arXiv e-print for this exact version is bundled unchanged as \texttt{sources/195581/source0.tex}.
\begin{definition}[Admissible curve set]\label{def:AC}
An \emph{admissible curve set} for $C$ is a finite collection $\mathcal A_C$ of isotopy classes of essential
simple closed curves and properly embedded essential arcs on $S_{g,n}$,
together with a designated involution
\[
\sigma_\ast:\mathcal A_C\longrightarrow\mathcal A_C,\qquad \gamma\longmapsto \gamma^\sigma,
\]
encoding which curve family on the $\mathrm{II}$--side is paired with a given curve family on the $\mathrm{I}$--side.
\end{definition}

For the analytic part of the paper, we fix an admissible curve set $\mathcal A_C$ and choose
continuous extremal--length assignments $\Ext_{\mathrm{I}},\,\Ext_{\mathrm{II}}$ in the sense of
Definition~\ref{def:ExtAssign}.
Concretely, for each curve family $\gamma\in\mathcal A_C$ we are given two positive functions
\[
u\longmapsto \Ext_{\mathrm{I}}(u;\gamma),\qquad
u\longmapsto \Ext_{\mathrm{II}}(u;\gamma)
\]
on $\Delta_C^{\mathrm{sc}}$. For $\gamma\in\mathcal A_C$, and \(u\in\Delta_C^{\mathrm{sc}}\) define the \emph{mismatch}
\[
m_\gamma(u)\ :=\ \log\frac{\Ext_{\mathrm{I}}(u;\gamma)}{\Ext_{\mathrm{II}}(u;\gamma^\sigma)},
\]
and  the \emph{height}
\begin{equation}\label{eq:height-def}
H(u)\ :=\ \sum_{\gamma\in\mathcal A_C} m_\gamma(u)^2.
\end{equation}
We call $u\in\Delta_C^{\mathrm{sc}}$ \emph{reflexive} if $H(u)=0$, equivalently if
\[
\Ext_{\mathrm{I}}(u;\gamma)=\Ext_{\mathrm{II}}(u;\gamma^\sigma)\qquad\text{for all }\gamma\in\mathcal A_C.
\]
\begin{Hypotheses}\label{def:axioms}
We say that $(C,\mathcal A_C,\Ext_{\mathrm{I}},\Ext_{\mathrm{II}})$ satisfies the a hypotheses if:
\begin{enumerate}[label=\textup{(H\arabic*)}, leftmargin=2.8em]
\item it is Teichm\"uller regular in the sense of Definition~\ref{def:TeichReg};
\item it is degeneration--detecting in the sense of Definition~\ref{def:DegDetect};
\item it is pushable in the sense of Definition~\ref{def:Pushable}.
\end{enumerate}
\end{Hypotheses}

\begin{theorem}[Existence of a reflexive point]\label{thm:main}
Let $C$ be a configuration datum and let $\Delta_C$ be the associated restricted character domain, with scale--fixed slice
$\Delta_C^{\mathrm{sc}}$.
Fix an admissible curve set $\mathcal A_C$ and continuous extremal--length assignments $\Ext_{\mathrm{I}},\Ext_{\mathrm{II}}$,
and define $H$ by \eqref{eq:height-def}.
If the  hypotheses~\ref{def:axioms} hold, then $\Delta_C^{\mathrm{sc}}$ contains a reflexive point.
\end{theorem}

\begin{definition}[Extremal--length assignment]\label{def:ExtAssign}
An \emph{extremal--length assignment} for $C$ consists of two maps
\[
\Ext_{\mathrm{I}}:\Delta_C^{\mathrm{sc}}\times \mathcal A_C\longrightarrow \R_{>0},
\qquad
\Ext_{\mathrm{II}}:\Delta_C^{\mathrm{sc}}\times \mathcal A_C\longrightarrow \R_{>0},
\]
such that for each fixed $\gamma\in\mathcal A_C$ the functions
\[
u\longmapsto \Ext_{\star}(u;\gamma),\qquad \star\in\{\mathrm{I},\mathrm{II}\},
\]
are continuous on $\Delta_C^{\mathrm{sc}}$.
\end{definition}

\begin{definition}[Teichm\"uller regularity]\label{def:TeichReg}
We say that $C$ is \emph{Teichm\"uller regular} with respect to the chosen curve set and extremal--length assignment if:
\begin{enumerate}[label=\textup{(R\arabic*)}, leftmargin=2.8em]
\item for each $\gamma\in\mathcal A_C$, the functions $u\longmapsto \Ext_{\star}(u;\gamma)$ are $C^1$ on $\Delta_C^{\mathrm{sc}}$ for
$\star\in\{\mathrm{I},\mathrm{II}\}$;
\item the log--extremal--length maps
\[
\mathbf{E}_{\star}:\Delta_C^{\mathrm{sc}}\longrightarrow \R^{\mathcal A_C},\qquad
\mathbf{E}_{\star}(u):=\bigl(\log \Ext_{\star}(u;\gamma)\bigr)_{\gamma\in\mathcal A_C},
\]
are local immersions for $\star\in\{\mathrm{I},\mathrm{II}\}$.
\end{enumerate}
\end{definition}

\begin{definition}[Degeneration detection]\label{def:DegDetect}
We say that $(C,\mathcal A_C)$ is \emph{degeneration--detecting} if for every sequence
$u_k\in\Delta_C^{\mathrm{sc}}$ leaving every compact subset of $\Delta_C^{\mathrm{sc}}$, there exists a curve
$\gamma\in\mathcal A_C$ such that
\[
|m_\gamma(u_k)|\longrightarrow\infty .
\]
\end{definition}

\begin{definition}[Pushability]\label{def:Pushable}
We say that $(C,\mathcal A_C)$ is \emph{pushable} if for each $\gamma\in\mathcal A_C$ there exist:
\begin{itemize}[leftmargin=2.2em]
\item a $C^1$ vector field $V_\gamma$ on $\Delta_C^{\mathrm{sc}}$, and
\item a finite incidence set $I(\gamma)\subset \mathcal A_C$ with $\gamma\in I(\gamma)$,
\end{itemize}
such that for every $u\in\Delta_C^{\mathrm{sc}}$ the following hold.
\begin{enumerate}[label=\textup{(P\arabic*)}, leftmargin=2.8em]
\item 
If $m_\gamma(u)\neq 0$, then
\(
\mathrm{sign}\bigl(m_\gamma(u)\bigr)\cdot (dm_\gamma)_u(V_\gamma) <0.
\)
\item 
If $\delta\notin I(\gamma)$ then $(dm_\delta)_u(V_\gamma)=0$. For $\delta\in I(\gamma)$ the derivatives
$(dm_\delta)_u(V_\gamma)$ are uniformly bounded in $u$.
\item 
If $m_\gamma(u)\neq 0$, then
\(
|m_\gamma(u)|\cdot\bigl|(dm_\gamma)_u(V_\gamma)\bigr|
\;>\;
\sum_{\delta\in I(\gamma)\setminus\{\gamma\}}
|m_\delta(u)|\cdot\bigl|(dm_\delta)_u(V_\gamma)\bigr|.
\)
\end{enumerate}
\end{definition}

\begin{proposition}[Core curves on \(X(1,b,c)\) with slit length \(\ell\)]\label{prop:ExtCorePlumbing}
On \(X(1,b,c)\) one has the exact equalities
\[
\Ext_{X(1,b,c)}(\alpha_1)=\frac{1}{\,b-\ell\,},\qquad
\Ext_{X(1,b,c)}(\beta_1)=b,
\]
and
\[
\Ext_{X(1,b,c)}(\alpha_2)=\frac{1}{\,c-\ell\,},\qquad
\Ext_{X(1,b,c)}(\beta_2)=c.
\]
\end{proposition}

\begin{proposition}\label{prop:mismatch_db}
For \(u=(b,c)\in\Delta^{\mathrm{sc}}_{}\) one has
\[
m_{\alpha_1}(u)=\log\frac{\Ext_{\mathrm{I}}(u;\alpha_1)}{\Ext_{\mathrm{II}}(u;\alpha_1)}
=\log\frac{\Ext_{X(u)}(\alpha_1)}{\Ext_{X(u)}(\beta_1)}
=\log\frac{(1/(b-\ell))}{b}
=-\log\bigl(b(b-\ell)\bigr),
\]
\[
m_{\alpha_2}(u)=\log\frac{\Ext_{\mathrm{I}}(u;\alpha_2)}{\Ext_{\mathrm{II}}(u;\alpha_2)}
=\log\frac{\Ext_{X(u)}(\alpha_2)}{\Ext_{X(u)}(\beta_2)}
=\log\frac{(1/(c-\ell))}{c}
=-\log\bigl(c(c-\ell)\bigr).
\]
Consequently,
\[
H(u)=m_{\alpha_1}(u)^2+m_{\alpha_2}(u)^2
=\bigl(\log(b(b-\ell))\bigr)^2+\bigl(\log(c(c-\ell))\bigr)^2.
\]
\end{proposition}

\begin{theorem}\label{thm:good_example}
With \(\mathcal A_{}\), \(\sigma_\ast=\id\), and the extremal--length assignments above, the
dumbbell domain \(\Delta^{\mathrm{sc}}_{}\cong\R_{>0}^2\) is Teichm\"uller regular
(Definition~\ref{def:TeichReg}), degeneration--detecting (Definition~\ref{def:DegDetect}), and pushable
(Definition~\ref{def:Pushable}).  Hence \(\Delta^{\mathrm{sc}}_{}\) contains a reflexive point.
In fact the unique reflexive point is \((b,c)=\bigl(b_\ell,b_\ell\bigr)\), where \(b_\ell> \ell\) is the unique positive root of \(b(b-\ell)=1\), i.e.\ \(b_\ell=\frac{\ell+\sqrt{\ell^2+4}}{2}\).
\end{theorem}

\begin{proof}
First we check: admissibility of \(\mathcal A_{}\).
The curves \(\alpha_1,\alpha_2\) are represented by simple closed geodesics on the translation surface
\(X(1,b,c)\), hence are essential and simple on \(S_2\).  Taking \(\sigma_\ast=\id\) gives an admissible curve
set in the sense of Definition~\ref{def:AC}.

\smallskip
\noindent

By Proposition~\ref{prop:ExtCorePlumbing},
\[
\log \Ext_{\mathrm{I}}(u;\alpha_1)=-\log(b-\ell),\qquad \log \Ext_{\mathrm{I}}(u;\alpha_2)=-\log(c-\ell),
\]
and
\[
\log \Ext_{\mathrm{II}}(u;\alpha_1)=\log b,\qquad \log \Ext_{\mathrm{II}}(u;\alpha_2)=\log c.
\]
Thus the extremal--length coordinate maps
\[
E_{\mathrm{I, II}}:\Delta^{\mathrm{sc}}_{}\longrightarrow \R^{\mathcal A_{}},
\qquad
E_{\mathrm{I}}(b,c)=(-\log(b-\ell),-\log(c-\ell)),\quad
E_{\mathrm{II}}(b,c)=(\log b,\log c)
\]
are \(C^\infty\) and have invertible Jacobian matrices \(\mathrm{diag}(-1/(b-\ell),-1/(c-\ell))\) and
\(\mathrm{diag}(1/b,1/c)\), respectively.
Hence they are local immersions, proving Teichm\"uller regularity.


\smallskip
\noindent

If \((b_k,c_k)\) leaves every compact subset of \(\Delta^{\mathrm{sc}}_{}=\{(b,c)\in\R_{>0}^2:\ \ell<\min\{b,c\}\}\),
then either \(b_k\longrightarrow \ell^+\) or \(b_k\longrightarrow\infty\), or \(c_k\longrightarrow \ell^+\) or \(c_k\longrightarrow\infty\).
By Proposition~\ref{prop:mismatch_db}, \(\bigl|\log(b_k(b_k-\ell))\bigr|\longrightarrow\infty\) or \(\bigl|\log(c_k(c_k-\ell))\bigr|\longrightarrow\infty\),
hence \(H(b_k,c_k)\longrightarrow\infty\).  Thus the height detects degeneration.

\smallskip
\noindent
Last thing to check pushability. Let
\[
V_{\alpha_1}:=b\,\frac{\partial}{\partial b},
\qquad
V_{\alpha_2}:=c\,\frac{\partial}{\partial c}.
\]
Since \(m_{\alpha_1}(b,c)=-\log\bigl(b(b-\ell)\bigr)=-\log b-\log(b-\ell)\), we compute
\[
(dm_{\alpha_1})_u(V_{\alpha_1})=b\,\frac{\partial}{\partial b}\Bigl(-\log b-\log(b-\ell)\Bigr)
= -\Bigl(1+\frac{b}{b-\ell}\Bigr)\ <\ 0
\quad\text{whenever }m_{\alpha_1}(u)\neq 0,
\]
and similarly
\[
(dm_{\alpha_2})_u(V_{\alpha_2})=
 -\Bigl(1+\frac{c}{c-\ell}\Bigr)\ <\ 0
\quad\text{whenever }m_{\alpha_2}(u)\neq 0.
\]
Moreover, the cross--derivatives vanish.  Thus the pushability conditions
(Definition~\ref{def:Pushable}) hold with incidence sets \(I(\alpha_1)=\{\alpha_1\}\) and \(I(\alpha_2)=\{\alpha_2\}\).

\smallskip
Having verified the hypotheses, Theorem~\ref{thm:main} implies the existence of a reflexive point.
In the present example we can identify it explicitly from the formula in Proposition~\ref{prop:mismatch_db}:
\[
H(b,c)=\bigl(\log(b(b-\ell))\bigr)^2+\bigl(\log(c(c-\ell))\bigr)^2.
\]
Thus \(H(b,c)=0\) if and only if \(b(b-\ell)=1\) and \(c(c-\ell)=1\), i.e.\ \(b=c=b_\ell=\frac{\ell+\sqrt{\ell^2+4}}{2}\).
\end{proof}

\end{document}
